\honest{Zombies: The Movie}{Eliezer Yudkowsky}{2008}

I write a film. Generals learn that New York has been overrun by zombies, and not the kind from the film of 28 days ago: these are philosophical zombies, who behave exactly as we do but are not conscious. Video of the city before and after shows crowds, restaurants and couples holding hands. ``It's worse than I imagined,'' says a colonel. \nb{Chalmers had written in \textsc{The Conscious Mind} (1996) that phenomenal zombies were unlikely to be popular in Hollywood, ``as there would be obvious problems with their depiction.'' This scene shows the problem.}

A scientist explains that the disease is an ``epiphenomenal virus.'' Asked how we know it exists, he answers: ``The same way we know we're conscious.'' \nb{The joke compresses the argument of ``Zombies! Zombies?''. Chalmers's reply, quoted in that post, is that having the experiences is what justifies the belief.} No cure works, because every placebo has an effect; a homeopath cannot be found, and the Taoists ``refuse to do anything.''

David Chalmers, one of the first victims, sits in a glass cell and says he feels fine, adding that he knows he would be expected to say that. He protests: ``This sort of thing can't actually happen!'' Asked why not, he gets as far as ``Because.'' \nb{The protest is accurate: Chalmers holds that zombies are unlikely to be naturally possible, because laws link physical states to experience in our world. The police badge in the next scene, ``Bridging Law Enforcement Agency,'' alludes to those laws.}

Daniel Dennett rides up to the prison on a motorcycle. Asked to show his qualia, he says he has none, which fits neither zombies nor humans, since both claim to have them. He disarms the guards, remarks ``I am one with my body,'' and says, ``There is a spoon.''

The zombies' one reply to the generals says that we are the ones with the virus. A colonel stares at his hands and says, ``I don't feel anything.'' The screen goes black, the sound goes silent, and ``The movie continues exactly as before.'' \nb{That line makes the case of the earlier posts in one sentence: removing something that has no effects changes nothing that happens.}

In a postscript I show a photograph of myself being attacked by zombie nurses at Penguicon, and tell a story about a knife and a question about the prior.
